% Este archivo es una biblioteca de derivaciones. Cada bloque se inserta en el
% lugar del cuerpo donde se usa, mediante \HAFOderivationpart. No se imprime
% como apéndice independiente.
\ifnum\HAFOderivationpart=1\relax
\subsection{Derivación completa de la forma de Lur'e}
\label{app:verificacion_lure}

Sea \(h(x)=\ell x+\psi(x)\). La primera ecuación de los dos modelos es
\[
{}^{\mathrm C}D_{t_0}^{q}x
=\alpha(y-x-h(x)).
\]
Al distribuir \(\alpha\),
\begin{align}
\alpha(y-x-h(x))
&=\alpha y-\alpha x-\alpha\ell x-\alpha\psi(x)\nonumber\\
&=-\alpha(1+\ell)x+\alpha y-\alpha\psi(x).
\label{eq:expansion_primera_fila}
\end{align}
Por otra parte,
\[
PX=
\begin{pmatrix}
-\alpha(1+\ell)x+\alpha y\\
x-y+z\\
-\beta y-\gamma z
\end{pmatrix},
\qquad
b\psi(r^{\T}X)=
\begin{pmatrix}
-\alpha\psi(x)\\0\\0
\end{pmatrix}.
\]
La suma reproduce \cref{eq:expansion_primera_fila} y las otras dos
ecuaciones. Para la saturación se toma
\(\ell=m_1\) y
\(\psi=(m_0-m_1)\operatorname{sat}\); para la arctangente,
\(\ell=a_1\) y \(\psi=a_2\arctan(\rho\,\cdot)\).

\fi
\ifnum\HAFOderivationpart=2\relax
\subsection{Derivación completa del determinante y la transferencia}
\label{app:transferencia_explicita}

Con \(\lambda=s^q\),
\begin{equation}
\lambda I-P=
\begin{pmatrix}
\lambda+\alpha(1+\ell)&-\alpha&0\\
-1&\lambda+1&-1\\
0&\beta&\lambda+\gamma
\end{pmatrix}.
\label{eq:matriz_resolvente}
\end{equation}
El menor de la posición \((1,1)\) es
\begin{align}
Q(\lambda)
&=
\det
\begin{pmatrix}
\lambda+1&-1\\
\beta&\lambda+\gamma
\end{pmatrix}\nonumber\\
&=(\lambda+1)(\lambda+\gamma)+\beta.
\label{eq:menor_Q}
\end{align}
Al desarrollar \(\det(\lambda I-P)\) por la primera fila,
\begin{align}
\Delta_\ell(\lambda)
&=\bigl[\lambda+\alpha(1+\ell)\bigr]Q(\lambda)
-(-\alpha)
\det
\begin{pmatrix}
-1&-1\\
0&\lambda+\gamma
\end{pmatrix}\nonumber\\
&=\bigl[\lambda+\alpha(1+\ell)\bigr]Q(\lambda)
-\alpha(\lambda+\gamma).
\label{eq:determinante_transferencia}
\end{align}

Para calcular la primera componente de
\((\lambda I-P)^{-1}b\), la regla de Cramer sustituye la primera columna de
\cref{eq:matriz_resolvente} por
\(b=(-\alpha,0,0)^{\T}\). El determinante del numerador es
\[
\det
\begin{pmatrix}
-\alpha&-\alpha&0\\
0&\lambda+1&-1\\
0&\beta&\lambda+\gamma
\end{pmatrix}
=-\alpha Q(\lambda).
\]
Como \(r^{\T}\) selecciona esa primera componente,
\begin{equation}
W_q(s)
=r^{\T}(s^qI-P)^{-1}b
=-\alpha\frac{Q(s^q)}{\Delta_\ell(s^q)}.
\label{eq:transferencia_cramer}
\end{equation}

La identidad
\[
P-s^qI=-(s^qI-P)
\]
implica
\[
r^{\T}(P-s^qI)^{-1}b=-W_q(s).
\]
Por ello, cambiar el orden de los términos dentro de la resolvente cambia
simultáneamente el signo de la transferencia y el signo del cierre. No cambia
la solución física si ambas modificaciones se realizan de manera coherente.

\fi
\ifnum\HAFOderivationpart=3\relax
\subsection{Derivación completa de los equilibrios}
\label{app:equilibrios}

En un equilibrio \(e=(x_e,y_e,z_e)^{\T}\), las dos últimas ecuaciones de
ambos sistemas dan
\[
0=x_e-y_e+z_e,\qquad
0=-\beta y_e-\gamma z_e.
\]
De la primera,
\[
z_e=y_e-x_e.
\]
Al sustituirla en la segunda,
\begin{align*}
0
&=-\beta y_e-\gamma(y_e-x_e)\\
&=-(\beta+\gamma)y_e+\gamma x_e.
\end{align*}
Si \(\beta+\gamma\neq0\),
\begin{equation}
y_e=\frac{\gamma}{\beta+\gamma}x_e,\qquad
z_e=-\frac{\beta}{\beta+\gamma}x_e.
\label{eq:equilibrios_yz}
\end{equation}
La primera ecuación de equilibrio exige
\[
0=y_e-x_e-h(x_e).
\]
Mediante \cref{eq:equilibrios_yz},
\begin{equation}
h(x_e)+\frac{\beta}{\beta+\gamma}x_e=0.
\label{eq:equilibrio_escalar_comun}
\end{equation}
Esta única ecuación escalar determina todos los equilibrios; las otras dos
coordenadas se recuperan después con \cref{eq:equilibrios_yz}.

\subsubsection{Equilibrios del Chua no suave}

En la región \(|x_e|<1\), \(h(x_e)=m_0x_e\). Entonces
\begin{equation}
\left(m_0+\frac{\beta}{\beta+\gamma}\right)x_e=0.
\label{eq:equilibrio_interior_ns}
\end{equation}
Cuando el coeficiente no es cero, la única raíz interior es \(x_e=0\).

Para \(x_e>1\),
\[
h(x_e)=m_1x_e+(m_0-m_1),
\]
y \cref{eq:equilibrio_escalar_comun} queda
\[
\left(m_1+\frac{\beta}{\beta+\gamma}\right)x_e+(m_0-m_1)=0.
\]
Por tanto,
\begin{equation}
x_\star=
\frac{m_1-m_0}
{m_1+\dfrac{\beta}{\beta+\gamma}}.
\label{eq:raiz_exterior_ns}
\end{equation}
Si \(x_\star>1\), existe el equilibrio positivo. Como la no linealidad es
impar, existe también la raíz \(-x_\star<-1\). Los tres equilibrios son
\begin{equation}
E_0=(0,0,0),\qquad
E_\pm=
\left(
\pm x_\star,\,
\pm\frac{\gamma}{\beta+\gamma}x_\star,\,
\mp\frac{\beta}{\beta+\gamma}x_\star
\right).
\label{eq:equilibrios_ns_simbolicos}
\end{equation}
La condición regional \(x_\star>1\) debe comprobarse; no basta con resolver la
ecuación lineal exterior.

\subsubsection{Equilibrios del Chua con arctangente}

Al usar
\(h(x)=a_1x+a_2\arctan(\rho x)\) en
\cref{eq:equilibrio_escalar_comun}, se obtiene
\begin{equation}
\left(a_1+\frac{\beta}{\beta+\gamma}\right)x_e
+a_2\arctan(\rho x_e)=0.
\label{eq:equilibrio_escalar_arctan}
\end{equation}
La expresión es impar, de modo que \(x_e=0\) siempre es raíz y toda raíz
positiva no nula tiene una compañera negativa. Las raíces no nulas se
determinan numéricamente y se sustituyen en \cref{eq:equilibrios_yz}. Esta
ecuación se evalúa con \(\rho=1\) para el caso bibliográfico y con el valor de
\(\rho\) de \cref{eq:parametros_c590} para el caso \texttt{c590}; no deben
intercambiarse las dos parametrizaciones.

\fi
\ifnum\HAFOderivationpart=4\relax
\subsection{Derivación completa de los jacobianos regionales}
\label{app:jacobianos}

Para la forma de Lur'e,
\[
F(X)=PX+b\psi(r^{\T}X).
\]
La regla de la cadena produce
\begin{equation}
J_e=P+b\,\psi'(r^{\T}e)r^{\T}.
\label{eq:jacobiano_lure}
\end{equation}
Como \(r^{\T}e=x_e\), solamente cambia la entrada \((1,1)\).

En el sistema no suave,
\[
\psi'_{\mathrm{ns}}(x)=
\begin{cases}
m_0-m_1,&|x|<1,\\
0,&|x|>1.
\end{cases}
\]
Por ello,
\begin{equation}
J_{\mathrm{in}}=
\begin{pmatrix}
-\alpha(1+m_0)&\alpha&0\\
1&-1&1\\
0&-\beta&-\gamma
\end{pmatrix},
\qquad
J_{\mathrm{out}}=
\begin{pmatrix}
-\alpha(1+m_1)&\alpha&0\\
1&-1&1\\
0&-\beta&-\gamma
\end{pmatrix}.
\label{eq:jacobianos_ns}
\end{equation}
En \(x=\pm1\) la derivada clásica no existe. Los equilibrios usados en este
reporte no se encuentran sobre esas superficies, por lo que se emplea el
jacobiano de la región que contiene a cada uno.

Para la arctangente,
\[
\psi'_{\mathrm a}(x)
=\frac{a_2\rho}{1+\rho^2x^2},
\]
y
\begin{equation}
J_e=
\begin{pmatrix}
\displaystyle
-\alpha\left(1+a_1+
\frac{a_2\rho}{1+\rho^2x_e^2}\right)
&\alpha&0\\[3mm]
1&-1&1\\
0&-\beta&-\gamma
\end{pmatrix}.
\label{eq:jacobiano_arctan}
\end{equation}
Los valores propios de
\cref{eq:jacobianos_ns,eq:jacobiano_arctan} se clasifican con
\cref{eq:matignon}; la estabilidad local no sustituye el sondeo de cuenca.

\fi
\ifnum\HAFOderivationpart=5\relax
\subsection{Derivación completa del polinomio característico}
\label{app:polinomio_caracteristico}

Sea \(g=h'(x_e)\) la pendiente total de la no linealidad en un equilibrio.
El jacobiano de cualquiera de las dos familias puede escribirse como
\begin{equation}
J(g)=
\begin{pmatrix}
-\alpha(1+g)&\alpha&0\\
1&-1&1\\
0&-\beta&-\gamma
\end{pmatrix}.
\label{eq:jacobiano_pendiente_total}
\end{equation}
Por tanto,
\[
\lambda I-J(g)=
\begin{pmatrix}
\lambda+\alpha(1+g)&-\alpha&0\\
-1&\lambda+1&-1\\
0&\beta&\lambda+\gamma
\end{pmatrix}.
\]
El desarrollo por la primera fila, igual al de
\cref{eq:determinante_transferencia}, produce
\begin{equation}
\chi_g(\lambda)
=\bigl[\lambda+\alpha(1+g)\bigr]
\bigl[(\lambda+1)(\lambda+\gamma)+\beta\bigr]
-\alpha(\lambda+\gamma).
\label{eq:polinomio_caracteristico_agrupado}
\end{equation}
Como
\[
(\lambda+1)(\lambda+\gamma)+\beta
=\lambda^2+(1+\gamma)\lambda+(\beta+\gamma),
\]
la multiplicación y agrupación por potencias de \(\lambda\) da
\begin{equation}
\begin{aligned}
\chi_g(\lambda)
={}&\lambda^3+
\bigl[1+\gamma+\alpha(1+g)\bigr]\lambda^2\\
&+\bigl[
\beta+\gamma+\alpha(1+g)(1+\gamma)-\alpha
\bigr]\lambda\\
&+\alpha\bigl[(1+g)(\beta+\gamma)-\gamma\bigr].
\end{aligned}
\label{eq:polinomio_caracteristico_expandido}
\end{equation}
En el sistema no suave se sustituye \(g=m_0\) en \(E_0\) y \(g=m_1\)
en \(E_\pm\). En el sistema con arctangente,
\[
g=a_1+\frac{a_2\rho}{1+\rho^2x_e^2}.
\]
Las raíces de \cref{eq:polinomio_caracteristico_expandido} son los valores
propios utilizados en \cref{eq:matignon}; así, la clasificación local se
obtiene a partir del mismo campo que define la integración.

\fi
\ifnum\HAFOderivationpart=6\relax
\subsection{Derivación completa para la saturación centrada}
\label{app:df_saturacion_centrada}

Sea
\[
\psi_{\mathrm{ns}}(\sigma)=\Delta\operatorname{sat}(\sigma),
\qquad \Delta=m_0-m_1.
\]
Si \(0<A\leq1\), no hay saturación y
\[
\psi_{\mathrm{ns}}(A\sin\theta)=\Delta A\sin\theta.
\]
Como \(\int_0^\pi\sin^2\theta\dd\theta=\pi/2\),
\[
N_{\mathrm{ns}}(A)
=\frac{2}{\pi A}\Delta A\frac{\pi}{2}
=\Delta.
\]

Para \(A>1\), defínase
\[
\theta_0=\arcsin\left(\frac1A\right),
\qquad 0<\theta_0<\frac{\pi}{2}.
\]
En \(0\leq\theta\leq\theta_0\) y
\(\pi-\theta_0\leq\theta\leq\pi\), la entrada permanece en la parte lineal;
en el intervalo central la saturación vale \(1\). Entonces
\begin{align}
N_{\mathrm{ns}}(A)
&=\frac{2\Delta}{\pi A}
\left[
2A\int_0^{\theta_0}\sin^2\theta\dd\theta
+\int_{\theta_0}^{\pi-\theta_0}\sin\theta\dd\theta
\right]\nonumber\\
&=\frac{2\Delta}{\pi A}
\left[
A\theta_0-\frac{A}{2}\sin(2\theta_0)
+2\cos\theta_0
\right]\nonumber\\
&=\frac{2\Delta}{\pi A}
\left[A\theta_0+\cos\theta_0\right].
\label{eq:df_sat_pasos}
\end{align}
En el último paso se usó
\[
\frac{A}{2}\sin(2\theta_0)
=A\sin\theta_0\cos\theta_0
=\cos\theta_0.
\]
Finalmente,
\[
\cos\theta_0
=\sqrt{1-\frac1{A^2}}
=\frac{\sqrt{A^2-1}}{A},
\]
y \cref{eq:df_sat_pasos} se convierte en
\[
N_{\mathrm{ns}}(A)
=\Delta\frac{2}{\pi}
\left[
\arcsin\left(\frac1A\right)
+\frac{\sqrt{A^2-1}}{A^2}
\right],
\]
que coincide con \cref{eq:df_saturacion}.

\fi
\ifnum\HAFOderivationpart=7\relax
\subsection{Derivación completa de la saturación sesgada}
\label{app:df_saturacion_sesgada}

Considérese
\[
u(\theta)=c+A\cos\theta,\qquad A>0.
\]
Para escribir una fórmula que incluya los casos parcialmente y completamente
saturados, se define
\[
[v]_{-1}^{1}=\max(-1,\min(1,v))
\]
y los ángulos
\begin{equation}
\theta_+
=\arccos\left[\frac{1-c}{A}\right]_{-1}^{1},
\qquad
\theta_-
=\arccos\left[\frac{-1-c}{A}\right]_{-1}^{1}.
\label{eq:angulos_saturacion_sesgada}
\end{equation}
Como \((1-c)/A>(-1-c)/A\) y \(\arccos\) es decreciente,
\(0\leq\theta_+\leq\theta_-\leq\pi\). En el intervalo
\([0,\pi]\),
\[
\operatorname{sat}(u(\theta))=
\begin{cases}
1,&0\leq\theta\leq\theta_+,\\
c+A\cos\theta,&\theta_+\leq\theta\leq\theta_-,\\
-1,&\theta_-\leq\theta\leq\pi.
\end{cases}
\]
La simetría respecto de \(2\pi\) permite integrar solamente en
\([0,\pi]\). Para \(\psi=\Delta\operatorname{sat}\),
\begin{align}
\Psi_0(c,A)
=\frac{\Delta}{\pi}\Big[
&\theta_+
+c(\theta_--\theta_+)
+A(\sin\theta_--\sin\theta_+)\nonumber\\
&-(\pi-\theta_-)
\Big].
\label{eq:psi0_saturacion_cerrada}
\end{align}

Para la componente fundamental se usan
\[
\int\cos\theta\dd\theta=\sin\theta,\qquad
\int\cos^2\theta\dd\theta
=\frac{\theta}{2}+\frac{\sin2\theta}{4}.
\]
La sustitución por intervalos da
\begin{align}
\Psi_1(c,A)
=\frac{2\Delta}{\pi}\Bigg[
&\sin\theta_++\sin\theta_-
+c(\sin\theta_--\sin\theta_+)\nonumber\\
&+\frac{A}{2}(\theta_--\theta_+)
+\frac{A}{4}
\bigl(\sin2\theta_--\sin2\theta_+\bigr)
\Bigg].
\label{eq:psi1_saturacion_cerrada}
\end{align}
Las operaciones de recorte en
\cref{eq:angulos_saturacion_sesgada} hacen que
\cref{eq:psi0_saturacion_cerrada,eq:psi1_saturacion_cerrada} sigan siendo
válidas cuando alguna de las tres regiones tiene longitud cero.

La ecuación de media puede reducirse a una igualdad escalar. De
\cref{eq:transferencia_cramer},
\begin{align}
W_q(0)
&=-\alpha
\frac{\beta+\gamma}
{\alpha(1+\ell)(\beta+\gamma)-\alpha\gamma}\nonumber\\
&=-\frac{\beta+\gamma}
{(1+\ell)(\beta+\gamma)-\gamma}.
\label{eq:ganancia_dc}
\end{align}
Por tanto,
\begin{equation}
c=W_q(0)\Psi_0(c,A),
\label{eq:cierre_dc_escalar}
\end{equation}
y el cierre armónico es
\[
1-W_q(\ii\omega)\frac{\Psi_1(c,A)}{A}=0.
\]
Estas tres ecuaciones reales ---una para la media, y las partes real e
imaginaria del cierre--- determinan \((c,A,\omega)\).

\fi
\ifnum\HAFOderivationpart=8\relax
\subsection{Derivación completa para la arctangente}
\label{app:df_arctan}

Sea
\[
\psi_{\mathrm a}(\sigma)=a_2\arctan(\rho\sigma),
\qquad
\xi=\rho A.
\]
Se define
\[
I(\xi)=
\int_0^\pi\arctan(\xi\sin\theta)\sin\theta\dd\theta.
\]
Como \(I(0)=0\), puede calcularse \(I\) integrando su derivada:
\begin{align}
I'(\xi)
&=\int_0^\pi
\frac{\sin^2\theta}{1+\xi^2\sin^2\theta}\dd\theta\nonumber\\
&=\frac1{\xi^2}
\left[
\pi-\int_0^\pi
\frac{\dd\theta}{1+\xi^2\sin^2\theta}
\right]\nonumber\\
&=\frac{\pi}{\xi^2}
\left(1-\frac1{\sqrt{1+\xi^2}}\right).
\label{eq:derivada_integral_arctan}
\end{align}
En la última igualdad se usó la integral estándar
\[
\int_0^\pi\frac{\dd\theta}{1+\xi^2\sin^2\theta}
=\frac{\pi}{\sqrt{1+\xi^2}}.
\]
Una primitiva que se anula en \(\xi=0\) es
\begin{equation}
I(\xi)
=\pi\frac{\sqrt{1+\xi^2}-1}{\xi}.
\label{eq:integral_arctan}
\end{equation}
Al sustituir \(\xi=\rho A\) en
\cref{eq:df_centrada},
\begin{align}
N_{\mathrm a}(A)
&=\frac{2a_2}{\pi A}I(\rho A)\nonumber\\
&=\frac{2a_2}{\rho A^2}
\left(\sqrt{1+\rho^2A^2}-1\right)\nonumber\\
&=\frac{2a_2\rho}
{1+\sqrt{1+\rho^2A^2}}.
\label{eq:df_arctan_derivada}
\end{align}

Si el cierre exige \(N_{\mathrm a}(A)=k\), la amplitud puede despejarse.
De la forma racionalizada,
\[
1+\sqrt{1+\rho^2A^2}=\frac{2a_2\rho}{k}.
\]
Al aislar la raíz, elevar al cuadrado y simplificar,
\begin{equation}
A^2
=\frac{4a_2(a_2\rho-k)}{k^2\rho}.
\label{eq:amplitud_arctan}
\end{equation}
La solución se acepta solamente si el lado derecho es positivo y la igualdad
previa a elevar al cuadrado tiene lado derecho no menor que \(2\). De manera
equivalente, para \(A>0\), \(k\) tiene el mismo signo que \(a_2\) y se
encuentra estrictamente entre \(0\) y \(a_2\rho\), con el orden de los
extremos ajustado al signo de \(a_2\).

\fi
\ifnum\HAFOderivationpart=9\relax
\subsection{Derivación completa de la identidad modal y su normalización}
\label{app:identidad_modal}

Sea \(P_0=P+bkr^{\T}\). Entonces
\[
\zeta I-P_0=(\zeta I-P)-bkr^{\T}.
\]
Si \(\zeta I-P\) es invertible, el lema del determinante matricial
\(\det(M+uv^{\T})=\det(M)(1+v^{\T}M^{-1}u)\), con
\[
M=\zeta I-P,\qquad u=-bk,\qquad v=r,
\]
produce
\begin{align}
\det(\zeta I-P_0)
&=\det(\zeta I-P)
\left[1-kr^{\T}(\zeta I-P)^{-1}b\right]\nonumber\\
&=\det(\zeta I-P)\bigl[1-kG(\zeta)\bigr],
\label{eq:lema_determinante_aplicado}
\end{align}
donde
\[
G(\zeta)=r^{\T}(\zeta I-P)^{-1}b,
\qquad W_q(s)=G(s^q).
\]
Por continuidad, la identidad polinómica se extiende también a los valores en
los que la resolvente de \(P\) es singular. Si
\[
1-kG(\zeta_0)=0,
\]
entonces \(\det(\zeta_0I-P_0)=0\), de modo que existe
\(v\neq0\) con \(P_0v=\zeta_0v\). Siempre que \(r^{\T}v\neq0\), se reemplaza
\(v\) por \(v/(r^{\T}v)\) y se obtiene \(r^{\T}v=1\). Ésta es la
normalización usada en \cref{eq:modo_normalizado,eq:semilla_modal}.

En la rama sesgada,
\[
X_1=\bigl[(\ii\omega_0)^qI-P\bigr]^{-1}b\Psi_1.
\]
Al premultiplicar por \(r^{\T}\),
\[
r^{\T}X_1
=W_q(\ii\omega_0)\Psi_1
=W_q(\ii\omega_0)N_1A.
\]
El cierre \(1-W_qN_1=0\) prueba algebraicamente que
\(r^{\T}X_1=A\).

\fi
\ifnum\HAFOderivationpart=10\relax
\subsection{Derivación completa de la homotopía afín}
\label{app:homotopia_afin}

La aproximación sesgada separa
\[
\psi(\sigma)
=\Psi_0+N_1(\sigma-c)
+\bigl[\psi(\sigma)-\Psi_0-N_1(\sigma-c)\bigr].
\]
Se definen
\begin{equation}
P_{\mathrm a}=P+N_1br^{\T},\qquad
d_{\mathrm a}=b(\Psi_0-N_1c).
\label{eq:parametros_homotopia_afin}
\end{equation}
La familia de campos es
\begin{equation}
\begin{aligned}
F_\varepsilon(X)
&=P_{\mathrm a}X+d_{\mathrm a}\\
&\quad+\varepsilon b
\left[
\psi(r^{\T}X)-\Psi_0-N_1(r^{\T}X-c)
\right],
\qquad 0\leq\varepsilon\leq1.
\end{aligned}
\label{eq:homotopia_afin}
\end{equation}

En \(\varepsilon=0\), la media \(\bar X\) de
\cref{eq:cierre_dc_sesgado} es un equilibrio del auxiliar:
\begin{align*}
F_0(\bar X)
&=P\bar X+bN_1c+b\Psi_0-bN_1c\\
&=P\bar X+b\Psi_0=0.
\end{align*}
Para una perturbación \(\xi=X-\bar X\),
\[
F_0(\bar X+\xi)=P_{\mathrm a}\xi,
\]
de modo que el auxiliar conserva tanto la media como el modo armónico.

En \(\varepsilon=1\), al agrupar términos,
\begin{align*}
F_1(X)
&=PX+bN_1r^{\T}X+b\Psi_0-bN_1c\\
&\quad+b\psi(r^{\T}X)-b\Psi_0-bN_1r^{\T}X+bN_1c\\
&=PX+b\psi(r^{\T}X).
\end{align*}
Así, el extremo final de \cref{eq:homotopia_afin} es exactamente el sistema
físico de \cref{eq:lure_comun}. Una homotopía centrada de la forma
\(P_0X+\varepsilon b[\psi(\sigma)-k\sigma]\) no conserva, en general, la
media \(c\neq0\); por eso no se utiliza para la rama sesgada.
\fi
